\chapter{Exercício 1}

Vamos discretizar uma das equações de Maxwell em 1 dimensão no vácuo, 
como por exemplo

\begin{equation}\label{eq:maxwell-1d-vacuo}
	\pd{E_y}{t} =  - \frac{1}{\epsilon} \pd{H_z}{z} 
\end{equation}

Usando diferenças centradas, cada uma das derivadas de (\ref{eq:maxwell-1d-vacuo}) 
pode ser aproximada por 

\[ \pd{E_y}{t} = \frac{E^{n+1} - E^n}{\Delta t} + O\left( \Delta t^2 \right) \quad , \quad 
	\pd{H_z}{z} = \frac{H_{i+1/2} - H_{i-1/2}}{\Delta x} + O\left( \Delta x^2 \right)
\]

\noindent onde as diferenças centradas resultam em erros de ordem 2. O erro $\tau_E$ é dado pela 
diferença entre os lados esquerdo e direito da Equação (\ref{eq:maxwell-1d-vacuo}) após substituir 
as aproximações, ou seja,

\begin{align*}
\tau_E
&=
\frac{E^{n+1} - E^n}{\Delta t} + O\left( \Delta t^2 \right) - \frac{1}{\epsilon} \left( \frac{H_{i+1/2} - H_{i-1/2}}{\Delta x} + O\left( \Delta x^2 \right) \right) 
\\[6pt]
&=
\left[ \frac{E^{n+1} - E^n}{\Delta t} -  \frac{1}{\epsilon} \frac{H_{i+1/2} - H_{i-1/2}}{\Delta x} \right] + O\left( \Delta t^2 \right) + \frac{1}{\epsilon} O\left( \Delta x^2 \right)
\\[6pt]
&=
\left[ \frac{E^{n+1} - E^n}{\Delta t} -  \frac{1}{\epsilon} \frac{H_{i+1/2} - H_{i-1/2}}{\Delta x} \right] + C \Delta t^2  + \frac{C}{\epsilon} \Delta x^2
\\[6pt]
&\text{Aplicando a Equação \ref{eq:maxwell-1d-vacuo} para substituir o termo entre colchetes, temos:}
\\[6pt]
\tau_E
&\leq
C \left( \Delta t^2 + \Delta x^2  \right)
\end{align*}


\chapter{Exercício 2}

\section{Solução Analítica}

O primeiro passo é entender a solução analítica, exibida na Figura \ref{fig:ex2-analitico}. 
Temos uma onda estacionária, com um campo elétrico $E_y$ que oscila no tempo e tem um perfil senoidal no espaço.

\begin{figure}[htb]
	\caption{\label{fig:ex2-analitico}Solução analítica do exercício 2.}
	\begin{center}
        \includegraphics[width=0.8\textwidth,trim=1 1 1 1,clip]{ex2-analitico.png}
	\end{center}
\end{figure}

\section{Solução Numérica}

Todo o código usado nessa solução encontra-se no Anexo \ref{cap:codigo_ex2}. 

A solução obtida para os diferentes valores de 
$\Delta x$ e $\Delta t$ é exibida na Figura \ref{fig:ex2-numerico}.
Vemos que a solução numérica é parecida com a solução analítica 
em todos os valores de discretização usados. 

\begin{figure}[htb]
	\caption{\label{fig:ex2-numerico}Solução numérica do exercício 2.}
	\begin{center}
        \includegraphics[width=\textwidth,trim=1 1 1 1,clip]{ex2-numerico.png}
	\end{center}
\end{figure}

\section{Medida de erro}

\textcolor{red}{Dúvida: confirmar na aula como calcular isso.}
